\documentclass[11pt, a4paper]{article}
\usepackage[T1]{fontenc}
\usepackage[utf8]{inputenc}
\usepackage{amsmath, amssymb, amsthm}
\usepackage{bm}
\usepackage{geometry}
\usepackage{booktabs}
\usepackage{hyperref}
\usepackage{xcolor}
\usepackage{parskip}

\geometry{margin=2.5cm}

\newcommand{\vv}[1]{\bm{#1}}
\newcommand{\nn}{\hat{\vv{n}}}
\newcommand{\cross}[2]{#1 \times #2}
\newcommand{\norm}[1]{\left\|#1\right\|}

\title{\textbf{BeamContactMapping: Mathematical Formulation}\\
       \large Cosserat Plugin -- SOFA Framework}
\author{Cosserat Plugin Documentation}
\date{}

\begin{document}
\maketitle
\tableofcontents

% ─────────────────────────────────────────────────────────────────────────────
\section{Overview}
% ─────────────────────────────────────────────────────────────────────────────

\texttt{BeamContactMapping} is a \texttt{Multi2Mapping} that receives the
contact-pair descriptors produced by \texttt{SphereSweptIntersectionMethod}
(SSIM) and converts them into concrete contact-point positions in world space.
It also computes the Jacobian relating frame velocities to contact-point
velocities, which is required by the constraint solver.

\textbf{Inputs}
\begin{itemize}
  \item Beam~1 rigid frames: $\{\vv{p}^{(1)}_k, \vv{R}^{(1)}_k\}_{k=0}^{N_1-1}$
        (output of \texttt{DiscreteCosseratMapping}).
  \item Beam~2 rigid frames: $\{\vv{p}^{(2)}_k, \vv{R}^{(2)}_k\}_{k=0}^{N_2-1}$.
  \item Per-contact section indices $\{i, j\}$ and curvilinear parameters
        $\{\alpha, \beta\} \in [0,1]^2$ from SSIM.
  \item Beam radii $r_1, r_2 > 0$.
\end{itemize}

\textbf{Outputs} (for $K$ active contact pairs)
\begin{itemize}
  \item $2K$ contact-point positions in $\mathbb{R}^3$:
        index $2k$ is the Beam~1 surface point $\vv{P}^c_A$,
        index $2k+1$ is the Beam~2 surface point $\vv{P}^c_B$.
\end{itemize}

% ─────────────────────────────────────────────────────────────────────────────
\section{SSIM Output: What the Mapping Receives}
% ─────────────────────────────────────────────────────────────────────────────

For contact pair $k$, SSIM provides the scalar parameters
$(\alpha, \beta) \in [0,1]^2$ that minimise the surface-to-surface gap.

\subsection{ALGO\_1 -- Segment-to-Segment}

$\{i,j\}$ are segment indices.
The closest point on Beam~1 is located at parameter $\alpha$ on the
segment $[\vv{p}^{(1)}_i,\, \vv{p}^{(1)}_{i+1}]$, and the closest
point on Beam~2 at parameter $\beta$ on $[\vv{p}^{(2)}_j,\, \vv{p}^{(2)}_{j+1}]$.

\subsection{ALGO\_2 -- Node-to-Segment}

$i$ is a node index (not a segment index), $j$ is still a segment index.
The Beam~1 contact is located exactly at node $\vv{p}^{(1)}_i$, so
$\alpha = 0$ by construction.
The Beam~2 contact is at parameter $\beta$ on
$[\vv{p}^{(2)}_j,\, \vv{p}^{(2)}_{j+1}]$.

% ─────────────────────────────────────────────────────────────────────────────
\section{Forward Map: Contact-Point Positions (\texttt{apply})}
% ─────────────────────────────────────────────────────────────────────────────

\subsection{Centreline Interpolation}

\textbf{ALGO\_1}
\begin{align}
  \vv{P}_A &= (1-\alpha)\,\vv{p}^{(1)}_i + \alpha\,\vv{p}^{(1)}_{i+1},
  \label{eq:PA_algo1}\\
  \vv{P}_B &= (1-\beta)\,\vv{p}^{(2)}_j + \beta\,\vv{p}^{(2)}_{j+1}.
  \label{eq:PB}
\end{align}

\textbf{ALGO\_2}
\begin{align}
  \vv{P}_A &= \vv{p}^{(1)}_i
  \quad (\text{no interpolation; }{\alpha = 0}),
  \label{eq:PA_algo2}\\
  \vv{P}_B &= (1-\beta)\,\vv{p}^{(2)}_j + \beta\,\vv{p}^{(2)}_{j+1}.
  \label{eq:PB_algo2}
\end{align}

\subsection{Contact Normal}

\begin{equation}
  \nn = \frac{\vv{P}_A - \vv{P}_B}{\norm{\vv{P}_A - \vv{P}_B}}.
  \label{eq:normal}
\end{equation}

The normal points from Beam~2 centreline toward Beam~1 centreline.
If $\norm{\vv{P}_A - \vv{P}_B} < \varepsilon$ (coincident centrelines),
a fall-back direction is taken from the local $Y$-axis of frame $i$ of Beam~1.

\subsection{Surface Contact Points}

\begin{align}
  \vv{P}^c_A &= \vv{P}_A - r_1\,\nn
  && \rightarrow \text{output}[2k],
  \label{eq:PcA}\\
  \vv{P}^c_B &= \vv{P}_B + r_2\,\nn
  && \rightarrow \text{output}[2k+1].
  \label{eq:PcB}
\end{align}

The gap (signed distance between the two surfaces along $\nn$) is
\begin{equation}
  g = \nn \cdot (\vv{P}^c_B - \vv{P}^c_A)
    = \norm{\vv{P}_A - \vv{P}_B} - (r_1 + r_2).
  \label{eq:gap}
\end{equation}

Contact is active when $g \leq 0$.

\subsection{Radial Moment Arms}

Define the moment arm of each contact point relative to its interpolated
centreline point:
\begin{equation}
  \vv{a}_A = \vv{P}^c_A - \vv{P}_A = -r_1\,\nn,
  \qquad
  \vv{a}_B = \vv{P}^c_B - \vv{P}_B = +r_2\,\nn.
  \label{eq:arms}
\end{equation}

A key property used in the Jacobian derivation (Section~\ref{sec:jacobian})
is that $\vv{a}_A$ and $\vv{a}_B$ are the \emph{same} vector for all
frames contributing to the same contact point, because the SLERP
linearisation collapses per-frame moment variations into a single radial term.

% ─────────────────────────────────────────────────────────────────────────────
\section{Tangent Map: Velocity Jacobian (\texttt{applyJ})}
\label{sec:jacobian}
% ─────────────────────────────────────────────────────────────────────────────

\subsection{Rigid-Frame Velocity Representation}

Each Cosserat beam frame carries a $6$-DOF velocity:
$\vv{d}_k = (\vv{v}_k,\, \bm{\omega}_k)$
where $\vv{v}_k \in \mathbb{R}^3$ is the linear velocity and
$\bm{\omega}_k \in \mathbb{R}^3$ is the angular velocity of frame $k$.

The velocity of any point $\vv{q}$ rigidly attached to frame $k$ is
\begin{equation}
  \dot{\vv{q}} = \vv{v}_k + \bm{\omega}_k \times (\vv{q} - \vv{p}_k).
  \label{eq:rigid_vel}
\end{equation}

\subsection{Velocity of the Centreline Contact Point}

Differentiating the linear interpolation \eqref{eq:PA_algo1}
and applying the SLERP linearisation (angular velocity interpolates
linearly to first order in $\alpha$):

\begin{equation}
  \dot{\vv{P}}_A = (1-\alpha)\,\vv{v}_i + \alpha\,\vv{v}_{i+1}
  \quad \text{(ALGO\_1)}.
  \label{eq:Pdot_A}
\end{equation}

For ALGO\_2 ($\alpha = 0$):
\begin{equation}
  \dot{\vv{P}}_A = \vv{v}_i.
  \label{eq:Pdot_A2}
\end{equation}

\subsection{Velocity of the Surface Contact Point}

Differentiating \eqref{eq:PcA} with the contact normal $\nn$ frozen
at the current configuration:

\begin{align}
  \dot{\vv{P}}^c_A
  &= \dot{\vv{P}}_A - r_1\,\dot{\nn}
  \;\approx\; \dot{\vv{P}}_A + r_1\,\bigl[(1-\alpha)\,\bm{\omega}_i
              + \alpha\,\bm{\omega}_{i+1}\bigr] \times \nn.
  \label{eq:PcAdot_intermediate}
\end{align}

Factoring by frame:
\begin{equation}
  \boxed{
  \dot{\vv{P}}^c_A
  = \sum_{b \in \mathcal{B}_1(k)}
    w_b\,\bigl[\vv{v}_b + \bm{\omega}_b \times \vv{a}_A\bigr],
  }
  \label{eq:PcAdot}
\end{equation}

where $\mathcal{B}_1(k)$ and the weights $w_b$ are defined in
Table~\ref{tab:blocks}, and $\vv{a}_A = -r_1\,\nn$.

Similarly for Beam~2:
\begin{equation}
  \boxed{
  \dot{\vv{P}}^c_B
  = \sum_{b \in \mathcal{B}_2(k)}
    w_b\,\bigl[\vv{v}_b + \bm{\omega}_b \times \vv{a}_B\bigr],
  }
  \label{eq:PcBdot}
\end{equation}

with $\mathcal{B}_2(k)$ and $\vv{a}_B = +r_2\,\nn$.

\begin{table}[h]
\centering
\caption{Jacobian block sets for contact pair $k$.}
\label{tab:blocks}
\begin{tabular}{lllll}
\toprule
Algorithm & Surface pt & Block set & Frame index & Weight \\
\midrule
ALGO\_1 & $\vv{P}^c_A$ & $\mathcal{B}_1(k)$ & $i$   & $(1-\alpha)$ \\
        &               &                     & $i+1$ & $\alpha$ \\
\addlinespace
ALGO\_2 & $\vv{P}^c_A$ & $\mathcal{B}_1(k)$ & $i$ only & $1$ \\
\midrule
Both    & $\vv{P}^c_B$ & $\mathcal{B}_2(k)$ & $j$   & $(1-\beta)$ \\
        &               &                     & $j+1$ & $\beta$ \\
\bottomrule
\end{tabular}
\end{table}

\subsection{Compact Jacobian Notation}

Stack the frame velocities of Beam~1 into $\dot{\vv{q}}_1 \in \mathbb{R}^{6N_1}$
and Beam~2 into $\dot{\vv{q}}_2 \in \mathbb{R}^{6N_2}$.
The contact-point velocity $\dot{\vv{s}}_k \in \mathbb{R}^3$ at output DOF $2k$
satisfies

\begin{equation}
  \dot{\vv{s}}_k = J^{(1)}_k\,\dot{\vv{q}}_1,
  \qquad
  J^{(1)}_k \in \mathbb{R}^{3 \times 6N_1},
  \label{eq:J1}
\end{equation}

where the only non-zero $3\times 6$ blocks of $J^{(1)}_k$ are
\begin{equation}
  \left[J^{(1)}_k\right]_{\cdot,b}
  = w_b\,\bigl[\,I_{3\times 3}\;\Big|\; -[\vv{a}_A]_\times\,\bigr],
  \qquad b \in \mathcal{B}_1(k),
  \label{eq:Jblock1}
\end{equation}

with $[\vv{a}_A]_\times$ the skew-symmetric matrix associated with
$\vv{a}_A$ (so that $[\vv{a}_A]_\times\,\bm{\omega} = \vv{a}_A \times \bm{\omega}$).

Analogously, for output DOF $2k+1$:
\begin{equation}
  \left[J^{(2)}_k\right]_{\cdot,b}
  = w_b\,\bigl[\,I_{3\times 3}\;\Big|\; -[\vv{a}_B]_\times\,\bigr],
  \qquad b \in \mathcal{B}_2(k).
  \label{eq:Jblock2}
\end{equation}

% ─────────────────────────────────────────────────────────────────────────────
\section{Transpose Map: Force Back-Propagation (\texttt{applyJT})}
\label{sec:applyJT}
% ─────────────────────────────────────────────────────────────────────────────

\subsection{Virtual Work Principle}

By the principle of virtual work, a contact force $\vv{F}$ applied at
$\vv{P}^c_A$ is dual to the velocity $\dot{\vv{P}}^c_A$:
\begin{equation}
  \delta W = \vv{F} \cdot \delta\vv{P}^c_A
           = \sum_{b \in \mathcal{B}_1(k)}
             w_b\,\vv{F} \cdot \bigl[\delta\vv{p}_b
             + \delta\bm{\theta}_b \times \vv{a}_A\bigr].
  \label{eq:virtualwork}
\end{equation}

Identifying the conjugate generalised forces:

\begin{equation}
  \boxed{
  \vv{f}_b \mathrel{+}= w_b\,\vv{F},
  \qquad
  \bm{\tau}_b \mathrel{+}= w_b\,(\vv{a}_A \times \vv{F}),
  \qquad b \in \mathcal{B}_1(k).
  }
  \label{eq:Jt1}
\end{equation}

For Beam~2, a force $\vv{F}'$ at $\vv{P}^c_B$ gives:
\begin{equation}
  \boxed{
  \vv{f}_b \mathrel{+}= w_b\,\vv{F}',
  \qquad
  \bm{\tau}_b \mathrel{+}= w_b\,(\vv{a}_B \times \vv{F}'),
  \qquad b \in \mathcal{B}_2(k).
  }
  \label{eq:Jt2}
\end{equation}

\subsection{Normal vs.\ Tangential Forces}

\paragraph{Normal contact force.}
$\vv{F} = \lambda\,\nn$ for some $\lambda \geq 0$:
\begin{equation}
  \bm{\tau}_b
  = w_b\,(\vv{a}_A \times \lambda\,\nn)
  = w_b\,(-r_1\,\nn) \times (\lambda\,\nn)
  = -w_b\,r_1\,\lambda\,(\nn \times \nn)
  = \vv{0}.
  \label{eq:normal_torque}
\end{equation}
A pure normal force produces \emph{no torque} on any frame.
Physically: the force is directed along the line connecting the
frame centre to the contact point, so its moment arm is zero.

\paragraph{Tangential (friction) force.}
Let $\hat{\vv{t}} \perp \nn$ be a tangential direction.
$\vv{F} = \lambda_t\,\hat{\vv{t}}$:
\begin{equation}
  \bm{\tau}_b
  = w_b\,(-r_1\,\nn) \times (\lambda_t\,\hat{\vv{t}})
  = -w_b\,r_1\,\lambda_t\,(\nn \times \hat{\vv{t}})
  \neq \vv{0}.
  \label{eq:tangential_torque}
\end{equation}
A tangential force generates a rolling/twisting torque on the frames.

% ─────────────────────────────────────────────────────────────────────────────
\section{Constraint Jacobian (\texttt{applyJT} MatrixDeriv)}
\label{sec:constraint}
% ─────────────────────────────────────────────────────────────────────────────

\subsection{Role in the Constraint Solver}

\texttt{GenericConstraintSolver} (under \texttt{FreeMotionAnimationLoop})
assembles a constraint Jacobian $H$ by calling \texttt{applyJT} on each
mapping in the scene graph. Each \emph{constraint} corresponds to one
scalar condition $g(\vv{q}) = 0$, evaluated along a direction
$\vv{d} \in \mathbb{R}^3$ at a specific contact point.

A row $r$ of the output constraint matrix encodes:
\begin{equation}
  [\text{OutMatrixDeriv}]_{r,\,\text{outDOF}}
  = \vv{d} \in \mathbb{R}^3.
  \label{eq:out_constraint_row}
\end{equation}

For the normal contact constraint, $\vv{d} = \nn$ and
$\text{outDOF} \in \{2k,\, 2k+1\}$ corresponds to one of the two
surface points of pair $k$.

\subsection{Mapping to Input Space}

The corresponding contribution to the input (frame) constraint matrix is:

For $\text{outDOF} = 2k$ (Beam-1 surface point, blocks $b \in \mathcal{B}_1(k)$):
\begin{equation}
  [\text{In1MatrixDeriv}]_{r,\, b.\text{frameIdx}}
  \mathrel{+}=
  \begin{pmatrix} w_b\,\vv{d} \\ w_b\,(\vv{a}_A \times \vv{d}) \end{pmatrix}
  \in \mathbb{R}^6.
  \label{eq:in1_constraint}
\end{equation}

For $\text{outDOF} = 2k+1$ (Beam-2 surface point, blocks $b \in \mathcal{B}_2(k)$):
\begin{equation}
  [\text{In2MatrixDeriv}]_{r,\, b.\text{frameIdx}}
  \mathrel{+}=
  \begin{pmatrix} w_b\,\vv{d} \\ w_b\,(\vv{a}_B \times \vv{d}) \end{pmatrix}
  \in \mathbb{R}^6.
  \label{eq:in2_constraint}
\end{equation}

\subsection{Normal Constraint Row (Explicit Form)}

The gap rate for contact $k$ along $\nn$ is
\begin{equation}
  \dot{g} = \nn \cdot (\dot{\vv{P}}^c_B - \dot{\vv{P}}^c_A).
  \label{eq:gdot}
\end{equation}

Substituting \eqref{eq:PcAdot}--\eqref{eq:PcBdot}:

\begin{align}
  \dot{g}
  &= \nn \cdot \Biggl[
    \sum_{b \in \mathcal{B}_2(k)} w_b\,(\vv{v}_b + \bm{\omega}_b \times \vv{a}_B)
    -
    \sum_{b \in \mathcal{B}_1(k)} w_b\,(\vv{v}_b + \bm{\omega}_b \times \vv{a}_A)
  \Biggr].
  \label{eq:gdot_expanded}
\end{align}

Using $\nn \cdot (\bm{\omega} \times (-r_1\nn)) = -r_1\,\nn\cdot(\bm{\omega}\times\nn) = 0$
(the scalar triple product $[\nn,\,\nn,\,\bm{\omega}] = 0$), the
angular-velocity contributions to the \emph{normal} constraint vanish:

\begin{equation}
  \dot{g}
  = \sum_{b \in \mathcal{B}_2(k)} w_b\,\nn \cdot \vv{v}_b
  - \sum_{b \in \mathcal{B}_1(k)} w_b\,\nn \cdot \vv{v}_b,
  \label{eq:gdot_normal_simplified}
\end{equation}

confirming that the $6$-DOF Jacobian blocks for the normal constraint
reduce to:
\begin{equation}
  \bigl[J^{(1)}_k\bigr]_{\text{normal},\, b}
  = \bigl[(1-\alpha)\,\nn^T\;\big|\;\vv{0}_{1\times 3}\bigr],
  \quad
  \bigl[J^{(2)}_k\bigr]_{\text{normal},\, b}
  = \bigl[(1-\beta)\,\nn^T\;\big|\;\vv{0}_{1\times 3}\bigr].
  \label{eq:normal_Jblocks}
\end{equation}

\subsection{Tangential Constraint Row (Explicit Form)}

For a tangential direction $\hat{\vv{t}} \perp \nn$, the angular-velocity
terms do not cancel. Using the cyclic identity
$\hat{\vv{t}} \cdot (\bm{\omega} \times \vv{a}) = \bm{\omega} \cdot (\vv{a} \times \hat{\vv{t}})$:

\begin{equation}
  \bigl[J^{(1)}_k\bigr]_{\text{tan},\, b}
  = w_b\,\bigl[\hat{\vv{t}}^T\;\big|\; (\vv{a}_A \times \hat{\vv{t}})^T\bigr]
  = w_b\,\bigl[\hat{\vv{t}}^T\;\big|\; (-r_1\,\nn \times \hat{\vv{t}})^T\bigr].
  \label{eq:tangential_Jblocks}
\end{equation}

The rotational block $-r_1\,(\nn \times \hat{\vv{t}})$ is the moment-arm
coupling that links frame rotation to tangential slip.

% ─────────────────────────────────────────────────────────────────────────────
\section{ALGO\_1 vs.\ ALGO\_2: Summary of Differences}
% ─────────────────────────────────────────────────────────────────────────────

\begin{table}[h]
\centering
\caption{Key differences between ALGO\_1 and ALGO\_2.}
\label{tab:algo_diff}
\begin{tabular}{lll}
\toprule
Property & ALGO\_1 (segment-to-segment) & ALGO\_2 (node-to-segment) \\
\midrule
Index $i$ & Beam-1 \emph{segment} index & Beam-1 \emph{node} index \\
Parameter $\alpha$ & $\in (0,1)$ in general & $= 0$ always \\
Beam-1 blocks per contact & 2 (frames $i$ and $i+1$) & 1 (frame $i$ only) \\
Weights on Beam-1 & $(1-\alpha)$, $\alpha$ & $1$ \\
Beam-2 blocks per contact & \multicolumn{2}{c}{2 (frames $j$ and $j+1$), weights $(1-\beta)$ and $\beta$} \\
Output positions per contact & \multicolumn{2}{c}{2 (Pc\_A at $2k$, Pc\_B at $2k+1$)} \\
\bottomrule
\end{tabular}
\end{table}

ALGO\_2 can be seen as the degenerate limit of ALGO\_1 with $\alpha = 0$:
the weight of frame $i+1$ vanishes, and the mapping needs only one Beam-1
Jacobian block. The same formulae \eqref{eq:PcAdot}--\eqref{eq:Jt2}
apply in both cases without special-casing -- only the block set
$\mathcal{B}_1(k)$ and the interpretation of the index $i$ differ.

% ─────────────────────────────────────────────────────────────────────────────
\section{Symbol Glossary}
% ─────────────────────────────────────────────────────────────────────────────

\begin{tabular}{ll}
\toprule
Symbol & Meaning \\
\midrule
$\vv{p}^{(m)}_k$ & Position of frame $k$ on Beam $m$ \\
$\vv{R}^{(m)}_k$ & Orientation (rotation matrix) of frame $k$ on Beam $m$ \\
$N_m$            & Number of frames on Beam $m$ \\
$i$              & Beam-1 segment index (ALGO\_1) or node index (ALGO\_2) \\
$j$              & Beam-2 segment index \\
$\alpha,\beta$   & Curvilinear parameters $\in [0,1]$ from SSIM \\
$r_1,r_2$        & Cross-section radii of Beam~1 and Beam~2 \\
$\vv{P}_A, \vv{P}_B$ & Interpolated centreline contact points \\
$\nn$            & Unit contact normal (Beam-2 $\to$ Beam-1) \\
$\vv{P}^c_A, \vv{P}^c_B$ & Surface contact points \\
$g$              & Signed gap ($= 0$ at first contact, $< 0$ in penetration) \\
$\vv{a}_A, \vv{a}_B$ & Radial moment arms: $-r_1\nn$ and $+r_2\nn$ \\
$\vv{v}_k, \bm{\omega}_k$ & Linear and angular velocities of frame $k$ \\
$w_b$            & Interpolation weight of Jacobian block $b$ \\
$\vv{F}$         & Contact force vector \\
$\vv{f}_b, \bm{\tau}_b$ & Generalised force and torque accumulated on frame $b$ \\
$K$              & Number of active contact pairs \\
$\varepsilon$    & Numerical zero threshold ($10^{-14}$) \\
\bottomrule
\end{tabular}

\end{document}
